% ch2_e §2.11–2.12 (书页 68–76 / p082–p090)
\section{声子–声子相互作用}

从形式上看，非谐项最重要的效应，是它破坏了各个晶格模的动力学独立性。例如，假设我们把代换 (2.8) 与 (2.10) 代入 (2.119)。对于 Bravais 格子，可以写出
\begin{align*}
\mathcal{V}^{(3)} &= \sum_{ll'l''}\mathsf{A}_{ll'l''}\colon \mathbf{u}_l\,\mathbf{u}_{l'}\,\mathbf{u}_{l''}\\
&= \sum_{ll'l''}\sum_{\mathbf{q}\mathbf{q}'\mathbf{q}''}\mathsf{A}_{ll'l''}\colon \mathbf{U}_{\mathbf{q}}\mathbf{U}_{\mathbf{q}'}\mathbf{U}_{\mathbf{q}''}\exp\bigl\{i(\mathbf{q}\cdot\mathbf{l}+\mathbf{q}'\cdot\mathbf{l}'+\mathbf{q}''\cdot\mathbf{l}'')\bigr\},
\tag{2.125}
\end{align*}
这里我们对三阶系数张量、以及它与各个晶格位移矢量的标积，都使用了非常紧凑的记号。

这一项本应计入运动方程之中。我们可以像 (2.6) 中那样应用平移不变性原理，来说明这类项的性质。(2.125) 中张量的系数只应依赖于格点 $\mathbf{l}$、$\mathbf{l}'$ 和 $\mathbf{l}''$ 的相对位置。于是
\begin{equation*}
\mathsf{A}_{ll'l''}=\mathsf{A}(\mathbf{h}',\mathbf{h}'')\quad\text{其中}\quad \mathbf{h}'=\mathbf{l}'-\mathbf{l},\ \mathbf{h}''=\mathbf{l}''-\mathbf{l}:
\tag{2.126}
\end{equation*}
把它代入 (2.125)，便得到
\begin{align*}
\mathcal{V}^{(3)} &= \sum_{l\mathbf{h}'\mathbf{h}''}\sum_{\mathbf{q}\mathbf{q}'\mathbf{q}''}\mathsf{A}(\mathbf{h}',\mathbf{h}'')\colon \mathbf{U}_{\mathbf{q}}\mathbf{U}_{\mathbf{q}'}\mathbf{U}_{\mathbf{q}''}\,e^{i(\mathbf{q}+\mathbf{q}'+\mathbf{q}'')\cdot\mathbf{l}}\,e^{i\mathbf{q}'\cdot\mathbf{h}'}e^{i\mathbf{q}''\cdot\mathbf{h}''}\\
&= \sum_{\mathbf{q}\mathbf{q}'\mathbf{q}''}\Bigl\{\sum_{\mathbf{h}'\mathbf{h}''}\mathsf{A}(\mathbf{h}',\mathbf{h}'')\,e^{i\mathbf{q}'\cdot\mathbf{h}'}e^{i\mathbf{q}''\cdot\mathbf{h}''}\Bigr\}\colon \mathbf{U}_{\mathbf{q}}\mathbf{U}_{\mathbf{q}'}\mathbf{U}_{\mathbf{q}''}\sum_{l}e^{i(\mathbf{q}+\mathbf{q}'+\mathbf{q}'')\cdot\mathbf{l}}\\
&= \sum_{\mathbf{q}\mathbf{q}'\mathbf{q}''}N\,\delta_{\mathbf{q}+\mathbf{q}'+\mathbf{q}'',\,\mathbf{g}}\,\mathsf{A}(\mathbf{q}',\mathbf{q}'')\colon \mathbf{U}_{\mathbf{q}}\mathbf{U}_{\mathbf{q}'}\mathbf{U}_{\mathbf{q}''}.
\tag{2.127}
\end{align*}
我们看到，模间耦合系数的张量就是非谐项张量的一个 Fourier 变换。尤其值得注意的是，对 $\mathbf{l}$ 的求和引入了一个形如 (2.92) 的因子，只有当
\begin{equation*}
\mathbf{q}+\mathbf{q}'+\mathbf{q}''=\mathbf{g},
\tag{2.128}
\end{equation*}
时它才不为零，这里 $\mathbf{g}$ 照例是倒格矢或零。于是，相互耦合的各个过程，其晶体动量就要满足一个条件。

我们可以这样来理解。假设我们把非谐项当作对晶格模运动的微扰，这些运动受方程 (2.11) 支配。通过一个正则变换，这些方程可以约化为一组独立简谐振子的运动方程。众所周知，这类方程在量子力学中用一个哈密顿量来表示，而这个哈密顿量又可以再变换成一个含有各模中量子的产生算符与湮没算符的表达式。

事实上，我们发现可以写出
\begin{equation*}
\mathbf{U}_{\mathbf{q}}=-i\left(\frac{\hbar}{2NM\nu_{\mathbf{q}}}\right)^{1/2}\mathbf{e}_{\mathbf{q}}\,(a_{\mathbf{q}}-a^{*}_{-\mathbf{q}}),
\tag{2.129}
\end{equation*}
其中 $\mathbf{e}_{\mathbf{q}}$ 是谱中某一支内波矢为 $\mathbf{q}$ 的模在偏振方向上的单位矢量，而 $a_{\mathbf{q}}$、$a^{*}_{\mathbf{q}}$ 是该模中一个声子的湮没算符与产生算符。当然，偏振矢量由本征值方程 (2.13) 对频率 $\nu_{\mathbf{q}}$ 的解来确定。

用这些算符来表示，由谐和项导出的哈密顿量约化成非常简单的形式
\begin{equation*}
\mathcal{H}_0=\sum_{\text{偏振}}\sum_{\mathbf{q}}\hbar\nu_{\mathbf{q}}\,(a_{\mathbf{q}}a^{*}_{\mathbf{q}}+a^{*}_{\mathbf{q}}a_{\mathbf{q}}),
\tag{2.130}
\end{equation*}
其本征态就是声子态，譬如 $|n_{\mathbf{q}}\rangle$，即第 $\mathbf{q}$ 个模中有 $n_{\mathbf{q}}$ 个量子。把 (2.129) 代入非谐项 (2.127)，就得到一个非常复杂的表达式，其中含有算符 $a_{\mathbf{q}}$、$a_{\mathbf{q}'}$ 等的乘积。例如，有一项包含（在挪动了一些正负号之后）
\begin{equation*}
\delta_{\mathbf{q}-\mathbf{q}'-\mathbf{q}'',\,\mathbf{g}}\,a^{*}_{\mathbf{q}}a_{\mathbf{q}'}a_{\mathbf{q}''}.
\tag{2.131}
\end{equation*}
当把这一项作为微扰作用到 $\mathcal{H}_0$ 的本征态上时，会引起这样的跃迁：在态 $\mathbf{q}$ 中产生一个声子，同时在态 $\mathbf{q}'$ 和 $\mathbf{q}''$ 中各消灭一个声子。相应的选择定则是
\begin{equation*}
\mathbf{q}'+\mathbf{q}''=\mathbf{q}-\mathbf{g}.
\tag{2.132}
\end{equation*}

这与式 (2.100) 相类似。若 $\mathbf{g}$ 为零，我们就说晶体动量守恒：动量 $\hbar\mathbf{q}'$ 和 $\hbar\mathbf{q}''$ 被消灭，凑出了动量 $\hbar\mathbf{q}$。当 $\mathbf{g}$ 不为零时，晶体动量不再守恒——但它只能改变 $\hbar\mathbf{g}$ 这么多。于是，\emph{声子–声子相互作用}可以像外束 X 射线的散射那样，用 N 过程和倒逆过程（U 过程）来描述。显然还有一些别的过程，其中一个声子被消灭、两个声子被产生——即 (2.131) 的共轭过程（译注：原文印作 (2.130)）。

可以证明，真实过程必须满足能量守恒；对 (2.132) 的情形，必须有
\begin{equation*}
\hbar\nu_{\mathbf{q}'}+\hbar\nu_{\mathbf{q}''}=\hbar\nu_{\mathbf{q}}.
\tag{2.133}
\end{equation*}
这一论证的依据是晶格模的时间依赖性。

显然，声子–声子相互作用实际的系数非常复杂：它们由非谐系数导出，还同偏振矢量等混杂在一起。不同偏振的模之间真实跃迁的规则也很复杂，因为这类跃迁要求能量条件与晶体动量条件同时得到满足。同样清楚的是，还可以有更高阶的声子–声子相互作用，例如 $n$ 个声子被消灭而 $n'$ 个声子被产生。这类项既可能来自晶格势能 Taylor 展开中的 $n+n'$ 阶项，也可能来自微扰级数中的更高阶项。它们必须满足类似 (2.128) 的选择定则——晶体动量的总改变必须是 $\hbar\mathbf{g}$，这里 $\mathbf{g}$ 为零或倒格矢。声子–声子相互作用过程在晶格热传导理论中十分重要，我们将在 §7.10 再回到这一点。

对于\emph{量子晶体}（quantum crystals）——特别是固态氦——上述晶格动力学的研究路线需要修正；在这类晶体中，原子的零点运动（zero-point motion）与晶格间距相当。由一串阶数越来越高的系数所定义、相继描述越来越复杂的声子相互作用过程的含时微扰级数，不再绝对收敛：我们甚至会发现，在平衡晶格间距处对势能求导而算得的二阶谐和力常数 (2.4) 取\emph{负}值。要修补这套代数，我们必须求助于\emph{多体理论}（many-body theory）的标准技巧之一。例如，我们可以走一走\emph{图形方法或簇展开}（diagrams or cluster expansions）的路子，在把微扰级数中各种无穷多项的组合求和之后，对所有相互作用系数进行重整化。或者，循着更直接的物理直觉，仿照 Hartree 方法的精神，为原子设立试探波函数，并对参数寻找变分条件。两种做法的结果大体相同。系统的行为相当正常：有声子谱、有声子–声子相互作用等等，只是有效力常数和更高阶的系数必须\emph{自洽地}（self-consistently）确定。于是，应把 (2.4) 换成大致如下的形式
\begin{equation*}
\mathsf{G}_{ll'}=\left\langle\left|\,\frac{\partial^2\mathcal{V}}{\partial\mathbf{u}_l\,\partial\mathbf{u}_{l'}}\,\right|\right\rangle,
\tag{2.134}
\end{equation*}
其中势能对原子位移的导数，要在某个与力常数的这些取值相一致的声子态 $|\;\rangle$ 中作为\emph{期望值}（expectation value）来求值。

\section{非完整格子的振动}

真实的晶体绝不可能是完美的，也不可能是无限延伸的；我们必须考虑缺陷和表面对晶格振动的影响。这个课题是没有尽头的，会通向诸如完全无序系统（如液体和玻璃）的动力学这类棘手的问题。不过，如果我们只限于讨论孤立的缺陷——例如质量不同的替位原子（substitutional atoms）——那么如今已能理解其中不少有趣的物理现象。

论证的要点通常被淹没在一大堆琐碎的细节之中。为了降低代数运算表面上的复杂性，让我们采用一种抽象的矩阵记法；例如，\emph{完美的}（Bravais）格子振动的一般方程在这种记法下便写成
\begin{equation*}
\mathsf{L}(\nu^2)\,u=0
\tag{2.135}
\end{equation*}
它代表诸如
\begin{equation*}
M\nu^2\mathbf{u}_l-\sum_{l'}\mathsf{G}_{ll'}\cdot\mathbf{u}_{l'}=0.
\tag{2.136}
\end{equation*}
这些方程本质上与 (2.5) 相同；列矩阵 $u$ 代表频率为 $\nu$ 的模中各位移 $\mathbf{u}_l$ 的全体，此时原子间相互作用力为 $\mathsf{G}_{ll'}$。动力学矩阵（dynamical matrix）$\mathsf{L}$ 具有晶格的平移对称性。按 §2.1 的论证，在每个声子频率 $\nu=\nu_{\mathbf{q}}$ 处 (2.135) 都有一个根：因为这时我们能为相应的声子态求得一个（归一化的）偏振矢量 $\mathbf{U}_{\mathbf{q}}$
\begin{equation*}
\mathbf{u}_l=\mathbf{U}_{\mathbf{q}}\,e^{i\mathbf{q}\cdot\mathbf{l}}
\tag{2.137}
\end{equation*}
而这只需解本征矢方程 (2.13) 即可。

现在假设我们改变了 $\mathbf{l}=0$ 处（其实是任选的一个格点）原子的质量，改变量为 $\delta M$，并且（或者）改变了该原子与晶体中近邻的相互作用，改变量为 $\delta\mathsf{G}_{ll'}$。我们的动力学方程 (2.136) 这时就应读作
\begin{equation*}
M\nu^2\mathbf{u}_l-\sum_{l'}\mathsf{G}_{ll'}\cdot\mathbf{u}_{l'}+\delta M\,\nu^2\mathbf{u}_0+\sum_{l'}\delta\mathsf{G}_{ll'}\cdot\mathbf{u}_{l'}=0,
\tag{2.138}
\end{equation*}
我们把它符号式地记为
\begin{equation*}
\mathsf{L}u+\delta\mathsf{L}u=0.
\tag{2.139}
\end{equation*}
这样，物理系统的这些改变就可以表示为原先完美晶体的动力学矩阵的一种微扰。

在线性链的特殊情形（参照 (2.15)），这些方程多少可以直接看出解来。但 Green 函数法（Green function method）适用面更广，而且也复杂不了多少。诀窍是把 (2.138) 改写成代数上等价的形式
\begin{equation*}
(\mathbf{1}+\mathsf{R}\,\delta\mathsf{L})\,u=0,
\tag{2.140}
\end{equation*}
其中预解式（resolvent）即 Green 函数 $\mathsf{R}$ 是动力学矩阵 $\mathsf{L}$ 的逆矩阵。只要 $\mathsf{R}$ 存在而且算得出来，我们就有了解决 (2.140) 所代表的线性方程的头绪。

预解式必须是 $\nu^2$ 的函数，满足抽象关系式
\begin{equation*}
\mathsf{R}(\nu^2)\,\mathsf{L}(\nu^2)=\mathbf{1}.
\tag{2.141}
\end{equation*}
很容易验证，本征矢 (2.137) 恰好给出这一问题的一组完全解：利用 (2.13) 诸解的标准正交性等性质，我们发现矩阵 $\mathsf{R}(\nu^2)$ 可以表示为
\begin{equation*}
\mathsf{R}_{ll'}(\nu^2)=\frac{1}{NM}\sum_{\mathbf{q}}\frac{\mathbf{U}^{*}_{\mathbf{q}}\mathbf{U}_{\mathbf{q}}}{\nu^2-\nu^2_{\mathbf{q}}}\,e^{i\mathbf{q}\cdot(\mathbf{l}-\mathbf{l}')}.
\tag{2.142}
\end{equation*}
这里我们同样出于简单，假定对所有波矢 $\mathbf{q}$ 的求和已扩展到把声子偏振的各种模式都包括进来。这个公式还使用了并矢（dyadic）约定，以顾及动力学矩阵的笛卡尔张量特征；不过，这仍然只是我们本应为例如一维链写下的那类公式的一个相当平凡的推广。

原则上，我们可以对 $(\mathbf{l}-\mathbf{l}')$ 的一切取值算出 (2.142)。实际上这并不必要，因为我们假定微扰 $\delta\mathsf{L}$ 高度局域（localized）在缺陷格点的邻域内。例如，考虑同位素杂质（isotopic impurity）的情形，此时 (2.138) 中只有质量发生了变化。把 (2.142) 代入 (2.140)，便得到一个只含 $\mathbf{u}_0$——即该格点上的矢量位移——的方程。其中只出现 $\mathsf{R}_{00}(\nu^2)$ 的求和，即
\begin{equation*}
\Bigl\{\mathbf{1}+\frac{\delta M}{M}\,\frac{\nu^2}{N}\sum_{\mathbf{q}}\frac{\mathbf{U}^{*}_{\mathbf{q}}\mathbf{U}_{\mathbf{q}}}{\nu^2-\nu^2_{\mathbf{q}}}\Bigr\}\cdot\mathbf{u}_0=0.
\tag{2.143}
\end{equation*}
笛卡尔张量 $\{\;\}$ 的 $3\times3$ 行列式为零，这一条件加在 $\nu^2$ 上，从而给出非完整晶格某一简正模的振动频率。更复杂的情形——几个格点上的力常数都有变化——显然会给出类似的方程，只是行列式的阶数更高。

遗憾的是，这些方程很难求解：即便是同位素替代这种初等情形，算起来也十分繁复。但为了理解究竟会发生什么，让我们撇开诸如 $\mathbf{u}_0$ 的矢量性这类几何特征，把 (2.143) 当作一个标量方程来处理。简正模的频率是如下形式方程的根
\begin{equation*}
1+\left(\frac{\delta M}{M}\right)\frac{1}{N}\sum_{\mathbf{q}}\frac{\nu^2}{\nu^2-\nu^2_{\mathbf{q}}}=0.
\tag{2.144}
\end{equation*}

%FIG{34}: 对含同位素杂质的晶格的模，式 (2.143) 的图解求法。

这些根可以用作图法求得（图 34）。我们要找的是函数
\begin{equation*}
f(\nu^2)\equiv\frac{1}{N}\sum_{\mathbf{q}}\frac{\nu^2}{\nu^2-\nu^2_{\mathbf{q}}}
\tag{2.145}
\end{equation*}
与水平线 $(-M/\delta M)$ 相交的那些点。

例如，设 $\delta M$ 为负。(2.144) 的每个根必定位于 $f(\nu^2)$ 的某个极点 $\nu^2_{\mathbf{q}}$ 的上方；受扰系统的各简正模在频率上与完美晶体的各简正模交错排布。由于 $\nu_{\mathbf{q}}$ 的取值构成一个稠密的\emph{带}，其间隔像 $1/N$ 那样趋于零，新解中的大多数与原有的解无从区分。但是再看一眼图 34 就会注意到，最高的那个根并不受这种约束，它可以离开带的顶端 $\nu_{\max}$ 一段有限的距离。当 $\nu$ 大于 $\nu_{\max}$ 时，求和 (2.145) 可以近似地用对谱密度（spectral density）(2.65) 的一个积分来表示，即
\begin{equation*}
\bar{f}(\nu^2)=\int^{\nu^{\max}}\frac{\nu^2\mathcal{D}(\nu')}{\nu^2-\nu'^2}\,d\nu'.
\tag{2.146}
\end{equation*}
在这个函数等于 $(-M/\delta M)$ 的频率处，就出现一个与缺陷相联系的特殊模。

这个特殊模必定是\emph{局域}的。这来自一条普遍定理：\emph{晶格的振荡运动在简正模频谱之外的频率处，其波数为虚数}。以 §2.2 的线性链模型为例。若把 $q=i\kappa$ 代入 (2.16)，我们便解出了 (2.15)，其频率由 (2.18) 的类比式给出，即
\begin{equation*}
\nu_{\kappa}=\sqrt{\left(\frac{\alpha}{M}\right)}\;2\sinh\left(\frac{\kappa a}{2}\right),
\tag{2.147}
\end{equation*}
对 $\nu_{\kappa}$ 的任何取值它都能得到满足。当然，作为完美晶格的简正模，这个解析解在物理上是不允许的，因为位移振幅 $u_l$ 会沿链像 $\exp(\pm\kappa l)$ 那样指数式地增长或衰减，从而不能满足循环边界条件。然而在非完整晶体中，可以把缺陷邻域的条件安排成使振动振幅向四面八方衰减，于是在样品远处的边界上变得可以忽略。从物理上讲，图 34 中的杂质原子比晶体中周围的原子来得轻，它作独立“Einstein”振荡时的固有频率会比较高。观察到的就是位于扩展声子模带之上的一个简正模，其中只有少数几个近邻原子随杂质协同运动（图 35(a)）。

缺陷位置上局域振动的其他种种例子都可以类似地理解。例如，如果我们把某个原子周围的弹簧常量加硬一些，就会观察到同样的效应。另一方面，比平常更重一些的原子具有较低的 Einstein 频率，因此我们应当能看到一个局域模从某个正常频带之\emph{下}脱离出来——例如从双原子晶体的光学支中脱离出来。

说来也怪，表面模（surface modes）的理论竟循着同样的路数。当然，我们可以从经典弹性理论中搬用关于连续介质自由表面上 Rayleigh 波与 Love 波的标准讨论，或许还要对表面平面内波矢的大小加上一个 Debye 限制。但这些波对应于那些局域在表面上、并朝晶体内部指数式衰减的振动模。虽然离散晶格中相应的问题看上去非常复杂，

%FIG{35}: (a) 轻杂质在声子带之上的局域模。(b) 重杂质在声子带之内的共振模。

但若把系统处理成一列晶格平面——它们从表面平面处的一个重大缺陷（例如劲度为零的弹簧常量）出发，沿一维向晶体内部延伸——还是能学到许多东西。这样一来，这类模应当出现在双原子晶体中、声子谱的光学支与声学支之间的带隙里。表面模可能对微小晶粒的比热有可检测的贡献，并且在某些电子器件中很重要。

局域杂质模的特征频率可以通过红外波段的光吸收直接观测（见 §8.4）。而且，有意思的是，如果杂质的固有振动频率落在某个声子带之\emph{内}，也会产生一条有所展宽的吸收线——例如在 (2.143) 的简单同位素情形中取 $\delta M$ 为正。这是\emph{共振}（resonance）的一个例子；在峰频率处，振动振幅在杂质格点上要比晶体其余部分大得多，但它并不随距离指数式地消衰，因此不是严格局域的（图 35(b)）。不出所料，在这个频率上杂质对入射声子的散射也非常强烈。

事实上，共振的理论只不过是局域态理论的一种解析延拓。共振频率是 (2.144) 的解，只是把其中的求和 $f(\nu^2)$ 换成 $\bar{f}(\nu^2)$ 的积分 (2.146) 的主部；而共振的\emph{宽度}则由同一积分的虚部给出，只不过现在要在 $\nu<\nu_{\max}$ 的情形下取值。这类关系的证明通常是在“球对称势中量子力学粒子的 Schrödinger 方程之解”的语境下给出的，但这一论证普遍适用于连续介质和规则晶格中的波。
